\subsubsection{Zonas ambientales, renovación y selección por estados (ENV-16)}
\label{sec:sd4-ambiente-env16}

ENV-16 conserva el balance de gas, humedad y estados de diseño. ENV-19, sección~\ref{sec:sd4-ambiente-env19}, actualiza su selección: ocho motores en cuatro trenes, ramales de 160 mm, colectores de 355 mm y arquitectura hidráulica exterior de tratamiento. Los diámetros de 125/250 mm y la pareja 022 utilizados a continuación describen el dimensionamiento anterior y su insuficiencia; la adquisición usa la selección posterior y mantiene sus brechas expresas. MON-16 conserva el monitoreo seleccionado, sin convertirse en control de seguridad de gas.

ENV-16 conserva la obra nueva de 10 por 7 m y 3 m libres, con seis volúmenes distintos: TI/climatización de 27,5 m$^2$/82,5 m$^3$, UPS A y B de 5,5 m$^2$/16,5 m$^3$ cada uno, bancos A y B de 12 m$^2$/36 m$^3$ cada uno y circulación de 7,5 m$^2$/22,5 m$^3$. Las áreas suman 70 m$^2$ y los volúmenes 210 m$^3$. Trabajo exterior de 6 m$^2$, respaldo apagado de 3 m$^2$, generación/combustible y bancadas exteriores siguen separados. La Figura~\ref{fig:plano-recinto-proyectado} conserva las huellas de bancos, UPS y gabinetes; no representa todavía equipos nuevos de tratamiento de aire de bancos. Esta frontera impide utilizar el área o retorno TI para justificar seguridad, temperatura o humedad de un recinto energético independiente \parencite[RT-06.03/13/14]{pucv-transversal}.

\textbf{Mapa constructivo de aire.} Cada banco recibe aire tratado exclusivamente por su reposición baja en el pasillo de servicio, sin toma desde UPS ni TI. Se fija para cada banco una entrada con área libre efectiva de al menos 0,15 m$^2$; a 900 m$^3$/h su velocidad media es 1,667 m/s. El tamaño exterior de la rejilla se calcula desde su fracción libre, no desde sus dimensiones nominales. Cada gabinete tiene dos captaciones altas, cuatro por recinto, de 225 m$^3$/h cada una en el estado de diseño. Cada ramal de proyecto tiene diámetro interior 125 mm; el colector, 250 mm. En ambos casos la velocidad de cálculo es 5,093 m/s. Dos circuitos independientes de extracción por banco, uno de servicio y uno de reserva, conservan cada uno 900 m$^3$/h efectivos con filtros sucios. El circuito detenido tiene cierre antirretorno; el cambio de ventilador no reduce el caudal requerido. Se adquieren cuatro conjuntos de extracción, ocho captaciones y cuatro colectores principales; caudal instalado redundante no equivale a renovar simultáneamente 1.800 m$^3$/h por banco.

Cada captación tiene dos salidas aislables hacia los dos colectores propios del banco: dieciséis ramales de 125 mm en ambos bancos, ocho activos en servicio normal. El control selecciona el colector sano y recibe su caudal agregado; no deja dos captaciones perdidas al retirar un ventilador. Los colectores salen por cubierta, con pasos sellados resistentes al fuego según el proyecto de incendio, descarga por encima de la cubierta y sin recirculación hacia admisiones, salidas o unidades exteriores. Se reserva separación horizontal propia de 3 m entre descargas de bancos y cualquier toma de aire, con verificación de viento y obstáculos; esta cota de proyecto no se presenta como distancia reglamentaria universal. Las rutas no atraviesan pasillos ni el frente de 1,20 m de las puertas. Las tuberías de condensado y suministro de agua se conducen por una galería hidráulica separada de los tableros y de las barras DC, con bandeja continua y pendiente hacia puntos supervisados. No se descarga agua sobre bancos, soleras eléctricas o salidas. El tratamiento exterior deberá tener envolvente dimensional seleccionada, mantenimiento y rutas representados en el plano antes de compra: una reserva de superficie sin variante no completa RT-06.03.

\textbf{Caudal independiente del ajuste no acreditado de 20 A.} La guía 93T publica un cargador configurable de hasta 60 A, sin individualizar los pasos para fijar 20 A. Por ello ENV-16 dimensiona la renovación para 60 A agregados por UPS, sin sumar dos veces las cadenas en paralelo. Con las 240 celdas en serie, los parámetros conservadores de gas y la ley de Faraday ya desarrollados en la memoria de estados, resultan 7,5978 m$^3$/h de hidrógeno por banco y, para 900 m$^3$/h, 0,8442\,\% bajo mezcla ideal. Para ambos bancos la renovación es 1.800 m$^3$/h. No se rebaja este caudal a 300 m$^3$/h mientras no exista ajuste autorizado e inhibición efectiva que limiten la corriente. La cifra conserva carácter de cálculo conservador de electrólisis, no de emisión medida ni de caudal impuesto por CSB \parencites[p.~98]{p7-eaton-93t-guia}{csbManualJunio2025}{nistCODATA2022}.

En 36 m$^3$, sin extracción y manteniendo esa emisión, la subida inicial sería $100(7,5978)/(36\times60)=0,35175$ puntos porcentuales por minuto. Desde 0,4\,\% hasta 1,0\,\% transcurrirían apenas 1,7058 minutos. La función exigida al enclavamiento local no depende de CloudWatch ni de una notificación humana: al perder caudal efectivo, señal válida de H$_2$ o alimentación de control debe retirar la autorización de recarga, conservando la descarga de la UPS y transfiriendo servicio por la rama sana conforme a la secuencia eléctrica. Se fija prealarma propia a 0,4\,\% e intervención crítica a 0,8\,\%; el estado normal exige permanecer bajo 1\,\%. La selección debe acreditar detector, rango, resolución, zona peligrosa, relés, medición de flujo y maniobra compatibles; estos umbrales no sustituyen esa selección. La guía acredita mando manual de cargador, no un contacto autónomo para esa función. Si se desarrolla apertura externa de la entrada AC al rectificador como solución, impedirá cargar pero pasará la rama a descarga, con maniobra y gestión explícita de autonomía por validar; ENV-16 no la presenta como interfaz automática ya seleccionada.

La extracción continúa durante descarga, flotación y purga posterior. No se asume emisión residual nula al cortar el cargador. El sistema de incendio conserva extracción ante prealarma y no sella ni descarga agente en el banco mientras H$_2$, caudal o su supervisión no permitan la maniobra aprobada. Una lectura por debajo de 0,4\,\% no demuestra por sí sola seguridad de diez minutos con compuertas cerradas. El análisis de retención de agente, evolución residual de H$_2$ y seccionamiento seguro se recibe conjuntamente con el diseño de incendio; un retardo arbitrario no resuelve ese conflicto.

\textbf{Selección térmica que cambia con el caudal.} Se mantienen exterior de proyecto 35\,$^{\circ}$C/80\,\% HR, presión de cálculo 101,325 kPa, consigna interior 24\,$^{\circ}$C/45\,\% HR y banda 23--25\,$^{\circ}$C/40--55\,\% HR sin condensación. A 1.800 m$^3$/h, la reposición añade 40,5187 kg/h de agua, 27,5752 kW latentes aproximados y 6,7320 kW sensibles aproximados. El balance entálpico completo es 35,0322 kW, 17,5161 kW por banco; se utiliza una sola vez, sin sumarle otra vez sus componentes. La extracción mínima media de agua por banco alcanza 20,25935 kg/h, antes de otros aportes. Las cifras son de obra nueva, no meteorología observada ni potencia eléctrica de un equipo de catálogo.

\begin{table}[htbp]
\centering\small
\caption{Estados ambientales que requieren selección conjunta ENV-16}\label{tab:sd4-env16-estados}
\begin{tabularx}{\linewidth}{L{27mm}X X}
\hline
\EncabezadoTabla{Estado} & \EncabezadoTabla{Zona y carga física} & \EncabezadoTabla{Condición N+1 y energía} \\ \hline
Red/flotación y carga mínima & TI: 2,710 kW máximos de equipos interiores, envolvente y control de vapor de hasta 0,50 kg/h; UPS: pérdidas de ambas unidades. Bancos: reposición separada; no usar calor de recarga simultáneo a descarga. & Una unidad de reserva por zona cubierta; comprobar mínimos ON/OFF, vapor y recalentamiento sin producir oscilación T/HR. Ventilación supervisada permanece activa. \\ \hline
Falla o mantenimiento de A/B & Parque completo en camino sano; UPS sana, consumo real del camino restante y extracción de ambos bancos según su estado. & No atribuir mitad de calor ni de potencia a una PSU por reparto supuesto. Fuente y capacidad de tratamiento del banco retirado siguen identificadas. \\ \hline
Descarga & Pérdidas DC/AC de UPS y calor neto de banco, con aire exterior tratado y gas residual. & Clima, extracción, detectores y control deben conservar energía durante la autonomía exigida; estar antes de UPS no otorga esa energía. \\ \hline
Retorno/recarga & Hasta 35,28 kW DC por UPS a 60 A/588 V; entrada de reposición de 17,5161 kW térmicos por banco, más calor neto efectivo de baterías y cargador. & No presentar los 35,28 kW DC como calor medido ni sumar recarga a la cota completa de rectificador. Selección térmica y eléctrica con las dos recuperaciones físicamente posibles. \\ \hline
\end{tabularx}
\FuenteTabla{Balance propio de Synaptix. Las pérdidas de 2,80092 kW de 93T son un punto AC/AC; no son un máximo certificado para todos los estados.}
\end{table}

La pareja candidata 022/PUZ-ZM250YKA2 entrega 22,6 kW totales y 19,3 kW sensibles brutos en el punto de ficha 27\,$^{\circ}$C/47\,\% HR de retorno y 35\,$^{\circ}$C exteriores; el ventilador de 0,70 kW deja 18,60 kW sensibles netos. Sus 3,3 kW latentes nominales equivalen aproximadamente a 4,849 kg/h usando 2.450 kJ/kg. Por banco, 20,25935/4,849=4,178 veces esa extracción nominal. Aun sin añadir calor de recarga, la pareja no acredita tratamiento de bancos; instalarla simplemente más grande en la misma sala TI tampoco corrige la separación de aire. La cota sensible previa de 16,7648 kW y los márgenes de fase siguen condicionados a la configuración anterior, no se extienden a ENV-16 \parencite[p.~3]{p6-mext-pi-2025}.

Los 60 A son techo configurable y escenario de dimensionamiento de extracción, no permiso para recargar así a cualquier carga de salida. La entrada del rectificador y sus pérdidas limitan la concurrencia: la memoria eléctrica conserva el máximo de entrada completo y exige autorización de corriente por estado de parque. La ventilación dimensionada para el extremo no obliga a reproducir simultáneamente extremos eléctricos físicamente incompatibles ni acredita recuperación a 60 A con parque pleno.

Los rangos de control se cotejan por variante: la guía 93T admite operación de 0 a 40\,$^{\circ}$C y 5--95\,\% HR sin condensación, con 25\,$^{\circ}$C recomendado; CSB individualiza las condiciones de carga/descarga y vida del bloque seleccionado. El Databook s-MEXT describe 19--35\,$^{\circ}$C de bulbo seco y 30--60\,\% HR, pero también 14--22,5\,$^{\circ}$C de bulbo húmedo y una envolvente gráfica; no basta comprobar dos intervalos independientes. Debe comprobarse cada esquina de la banda en el manual aplicable a la pareja YKA2, con exterior, tubería y accesorios seleccionados. Tampoco se trasladan curvas de YKA sin sufijo 2. Los límites particulares de servidores, comunicaciones, control y sensores se incorporan a esa intersección antes de compra \parencites[secc.~8]{p7-eaton-93t-guia}{csbHRL12540WRA260131}[p.~16]{p6-mext-databook}.

Para carga mínima no se presupone capacidad proporcional a los pasos de demanda. ME38 mantiene tiempos mínimos y límites específicos de deshumidificación; la selección exigirá mapa sensible/latente a carga baja, vapor de reposición y calor real por estado. Se conserva agua potable dentro de la tabla de calidad del humidificador 4301, válvulas de corte, bandejas y drenaje individual con alarma de bloqueo; no se sustituye por agua desmineralizada. Cada circuito R32 verifica carga total instalada en su volumen real, accesorios de seguridad, holguras, ventilación permanente, tubería y categoría de acceso. El mínimo de 21 m$^2$ de una 022 F2 no es permiso para agregar circuitos o introducirlos en bancos de 12 m$^2$. Protección anticorrosiva, admisiones limpias y mantenimiento se reciben en variante autorizada para entorno marino \parencites[pp.~139--149]{p6-mext-control-me38}[p.~42]{p6-mext-databook}[secc.~2.2.5--2.2.7]{p9-mext-manual-2024}.

La compra ambiental queda condicionada a una selección de tratamiento de bancos con curvas, potencia/corriente por fase, variante de ventilador apta para H$_2$ y presión real de ductos, detector y enclavamiento aprobados. La documentación oficial confirma que existen familias de extracción para hidrógeno, pero una familia no identifica la variante ni su punto caudal/presión. ENV-16 no convierte la sensibilidad con COP 3 o ventiladores de 0,600 kW en catálogo. Su efecto eléctrico/anual se incorpora a ENE junto con fuente de batería y recuperación; RT-06.13 y la huella ambiental final de RT-06.03 permanecen por completar \parencite{lote16-systemair-ex}.

\subsubsection{Monitoreo ambiental seleccionado y recepción (MON-16)}
\label{sec:sd4-monitoreo-mon16}

MON-16 selecciona dos APC NetBotz Rack Monitor 250 con NMC3 \texttt{NBRK0250A}, dos \texttt{NBPD0150}, doce sensores T/HR \texttt{AP9335TH}, doce cables de agua \texttt{NBES0308}, diez extensiones \texttt{NBES0309} y dos balizas \texttt{AP9324}. No se emplea \texttt{NBES0309} como detector autónomo ni se usa P161 como detección de agua. Cada monitor incorpora seis puertos universales y cada pod otros seis; las doce sondas de cada conjunto ocupan esos doce puertos, sin excederlos. Los dos AP9335TH incluidos con los monitores forman parte de los doce: se adquieren diez adicionales. Se dispone además de una unidad de monitor, un pod, dos T/HR, dos cables de agua y dos extensiones apagados en reserva local. El manual de la variante 250A acredita expresamente AP9335TH, NBES0308 y NBPD0150 y conexión A-Link del pod; A-Link no se conecta a Ethernet \parencites[pp.~24--26]{lote16-netbotz-guia}{lote16-netbotz-instalacion}{lote16-netbotz-pod}{lote16-netbotz-agua}.

Cada conjunto se alimenta de una rama UPS distinta, A en L2 y B en L3; la toma conmutada del monitor permanece vacía. La guía fija consumo de entrada propio de 0,25 A separado de hasta 10 A de la salida conmutada. Se reserva 57,5 VA por conjunto a 230 V y 115 VA para los dos, con pod y baliza dentro del presupuesto recibido, sin sumar 10 A de una toma vacía. Para el balance activo se usa la misma cota de 57,5 W por conjunto como límite conservador, no consumo medido. Cada equipo usa una conexión Ethernet cableada propia hacia un switch distinto de gestión; las sondas usan cable, no radio. La pérdida de A deja al conjunto B alimentado y observando todas las zonas, y viceversa; los sensores no se presentan como dispositivos de doble alimentación.

\begin{table}[htbp]
\centering\small
\caption{Mapa de puertos por cada conjunto A/B de MON-16}\label{tab:sd4-mon16-puntos}
\begin{tabularx}{\linewidth}{L{22mm}X X}
\hline
\EncabezadoTabla{Puerto} & \EncabezadoTabla{Sonda T/HR en monitor} & \EncabezadoTabla{Agua en puerto homólogo del pod} \\ \hline
1 & TI: A en entrada superior de cómputo, B en entrada superior de comunicaciones; ambos antes de mezclar con retorno. & Galería hidráulica TI: bandejas de ambos climatizadores, alimentación de humidificadores y ruta de sus drenajes; 12,2 m por conjunto. \\ \hline
2 & UPS A: A en entrada de aire y B en retorno del compartimento, sin colocarlos contra disipador. & Suelo bajo ruta de tuberías y perímetro útil UPS A, separado de bornes y pasillo; 6,1 m por conjunto. \\ \hline
3 & UPS B: A en entrada de aire y B en retorno del compartimento. & Suelo y perímetro útil UPS B; 6,1 m por conjunto. \\ \hline
4 & Banco A: dos puntos a 1,2 y 1,8 m en el pasillo, junto a cadenas A1/A2, fuera de bolsillos altos de H$_2$. & Ingreso de reposición, ruta hidráulica, bajos de gabinetes y suelo banco A; 18,3 m por conjunto. \\ \hline
5 & Banco B: disposición reflejada, junto a cadenas B1/B2. & Ruta correspondiente de banco B; 18,3 m por conjunto. \\ \hline
6 & Respaldo apagado: sensores separados en espacio de stock, lejos de pared exterior. & Perímetro de respaldo y umbral de ingreso, 6,1 m por conjunto. \\ \hline
\end{tabularx}
\FuenteTabla{Cada fila se ejecuta en A y B con cable independiente. Las longitudes suman 67,1 m de cable detector por conjunto y 134,2 m en ambos: seis bases de 6,1 m más cinco extensiones de 6,1 m por conjunto. Cada ruta permanece bajo el máximo de 30,5 m del manual.}
\end{table}

Las rutas se fijan sin cortar cable detector; el sobrante continúa en serpentina dentro de bandejas/perímetro sin provocar cruces de puertas ni contacto con conductores. Los cables de sondas, A-Link y Ethernet tienen identificadores distintos, plano de conexión y límites de longitud del fabricante; se fijan tendidos de sensor de hasta 20 m y A-Link de hasta 20 m por conjunto. Los sensores ordinarios no se colocan dentro de áreas clasificadas de emisión de H$_2$ ni sustituyen detectores de gas. El proyecto sitúa las captaciones y gas en los compartimentos altos/gabinetes y mantiene los puntos MON en área no clasificada; la recepción confirma esa frontera y sustituye por instrumentación apta si la clasificación se extiende, antes de habilitar. La detección de agua no constituye aprobación de mantener tuberías sobre equipos eléctricos.

Se configuran consigna 24\,$^{\circ}$C/45\,\% HR, advertencia fuera de 23--25\,$^{\circ}$C o 40--55\,\% HR durante 60 s, alarma crítica a menos de 20 o más de 27\,$^{\circ}$C, menos de 35 o más de 60\,\% HR durante 30 s y agua en cualquier ruta sin retardo intencional. El retorno de UPS se interpreta por ubicación y no reemplaza la consigna de entrada. El stock adopta advertencia a 18--27\,$^{\circ}$C/35--60\,\% HR y alarma a menos de 15 o más de 30\,$^{\circ}$C o más de 65\,\% HR. Son umbrales propios de intervención temprana, no límites de catálogo. Las alarmas no se limpian hasta permanecer dos minutos dentro de banda; agua requiere inspección y rearme autorizado. Las sondas AP9335TH tienen operación publicada de 0--55\,$^{\circ}$C y 0--95\,\% HR; el monitor admite 0--45\,$^{\circ}$C y 10--95\,\% HR no condensante. La recepción coteja error/ubicación con patrones; si impide discriminar una banda estrecha, se usa margen de medida conservador y se calibra o cambia la sonda \parencites{lote16-netbotz-th}[p.~167]{lote16-netbotz-guia}.

\textbf{Entrega a observabilidad y pérdida de señal.} Se habilita SNMPv3 con autenticación y privacidad, lista de acceso restringida a los dos recolectores locales y credenciales propias; SNMPv1, HTTP y cuentas predeterminadas quedan deshabilitados. Los dos Dell R360 existentes consultan ambos monitores cada 15 s y reciben traps y Syslog; se carga la versión de PowerNet MIB aplicable, conserva su hash y se realiza prueba de cada OID real. Se reserva por R360 0,10 núcleo, 256 MiB y 8 GiB de cola persistente dentro del presupuesto local, sin añadir otro servidor. Una lectura conserva ID de zona/sonda, unidad, valor, hora de adquisición y recepción, estado de comunicación, calibración y umbrales; no se importa una lectura anterior como valor actual. Dos consultas consecutivas fallidas o estado desconectado emiten pérdida de señal en como máximo 30 s más timeout máximo de 5 s. Se prueba desconexión de cada sonda y pod, no sólo silencio de un trap \parencites[pp.~54, 82--85]{lote16-netbotz-guia}{lote16-netbotz-mib}.

Se generan 36 magnitudes: 24 de T/HR y 12 estados de agua. A cuatro capturas por minuto, un registro propio de 256 bytes por magnitud ocupa $36(4)(60)(24)(256)=53.084.160$ bytes/día. Catorce días consumen 708,75 MiB; la reserva de 8 GiB por recolector incluye índices, eventos y recuperación con holgura, y la capacidad se comprueba con alarmas simultáneas. Los agentes capturan ambos, pero un líder con concesión local publica; el cambio de líder usa claves \texttt{recinto/sonda/secuencia/tipo} para no duplicar notificaciones. La pérdida de WAN mantiene balizas, alarmas y registros locales. Al recuperar, se entrega historia con hora original y calidad de conexión; métricas atrasadas que exceden el plazo admitido por CloudWatch se conservan como eventos históricos, sin falsear hora ni estado presente.

El exportador publica en espacio \texttt{Synaptix/Recinto}, dimensiones \texttt{Sitio,Zona,Sensor}, y CloudWatch Logs almacena alarmas, pérdida/restitución de señal y acciones. Se agrupan 144 registros/minuto para salida propia de hasta 2 kB/s con margen; una prueba de drenaje limita inicialmente 200 registros/s para no desplazar eventos operacionales. El NOC atiende las alertas según severidad de SD5, con asignación y acuse; la pérdida de WAN se distingue de falta de sensor y no promete aviso remoto por un canal caído. Las balizas se activan localmente desde el monitor, aun con ambos recolectores o la nube indisponibles; la alarma por fallo total del monitor la genera el otro conjunto y el recolector superviviente. Ninguna cadena MON gobierna por software cloud la recarga, fuego o parada de UPS.

La aceptación inyecta calor/humedad controlados, moja cada ruta detectora con método permitido por fabricante, bloquea y restituye drenaje, desconecta sonda/pod/Ethernet/A/B y corta WAN. Comprueba presencia y ubicación de 24 dispositivos, las 36 magnitudes, umbrales/retardos, baliza local, señal perdida, batería, sello de tiempo, cola de catorce días y entrega histórica sin alertas duplicadas. Se entrega plano de puntos, configuración, MIB/OIDs, calibración, mantenimiento y resultado firmado. Son actuaciones futuras de recepción de la configuración elegida; MON-16 desarrolla documentalmente RT-06.14 sin afirmar instalación ni ensayo ya ejecutados.
